% ============================================================
% Chapter 2: Module Attributes
% ============================================================
\section{模块属性（Module Attributes）}

\subsection{基本属性}

\begin{figure}[H]
\centering
\includegraphics[width=0.85\textwidth]{image3}
\caption{Module Attributes 概览}
\end{figure}

\begin{figure}[H]
\centering
\includegraphics[width=0.85\textwidth]{image4}
\caption{Module Attributes 详解}
\end{figure}

参考：
\begin{itemize}[leftmargin=1.5em]
    \item \href{https://www.erlang.org/doc/reference_manual/modules}{Erlang -- Modules}
    \item \href{https://www.tutorialspoint.com/erlang/erlang_modules.htm}{Erlang - Modules (TutorialsPoint)}
\end{itemize}

\subsection{Tuple 形式的模块属性}

Tuple 形式的属性与普通属性没有任何差别，例如：

\begin{lstlisting}
-module(rabbit).
-author("TutorialPoint").
-version("1.0").
-module_description("This is an example module.").

-rabbit_boot_step({pre_boot,
    [{description, "rabbit boot start"}]}).

-rabbit_boot_step({codec_correctness_check,
    [{description, "codec correctness check"},
     {mfa, {rabbit_binary_generator,
            check_empty_frame_size, []}},
     {requires, pre_boot},
     {enables, external_infrastructure}]}).

-export([start/0]).
-export([start1/0]).

start() ->
    io:fwrite("Hello World\n"),
    ModuleInfo = rabbit:module_info(),
    Attributes = rabbit:module_info(attributes),
    io:format("~p~n", [Attributes]),
    finish.
\end{lstlisting}

\begin{notebox}[要点]
Any module attribute can be specified. The attributes are stored in the compiled code and can be retrieved by calling \lstinline|Module:module_info(attributes)|, or by using the module \lstinline|beam_lib(3)| in STDLIB.
Tuple 形式的属性值是叠加的。
\end{notebox}
